\documentclass{../base/canon}

%% canon_variables.tex — Valores por defecto de metadatos institucionales
%%
%% Incluir ANTES del \begin{document} para sobreescribir los defaults de la clase.
%% El template puede sobreescribir cualquier valor después de este \input.
%%
%% Uso en template:
%%   %% canon_variables.tex — Valores por defecto de metadatos institucionales
%%
%% Incluir ANTES del \begin{document} para sobreescribir los defaults de la clase.
%% El template puede sobreescribir cualquier valor después de este \input.
%%
%% Uso en template:
%%   %% canon_variables.tex — Valores por defecto de metadatos institucionales
%%
%% Incluir ANTES del \begin{document} para sobreescribir los defaults de la clase.
%% El template puede sobreescribir cualquier valor después de este \input.
%%
%% Uso en template:
%%   \input{../base/canon_variables}
%%   \SetDocTitle{Mi Informe Específico}   % sobreescribe el default

\SetDocLogo{../assets/logo.pdf}
\SetDocUnit{Tecnología \& Sistemas}
\SetContactUnit{Mesa de soporte}
\SetContactMembers{%
  Equipo CoreX & Soporte & soporte@empresa.com%
}
\SetContactNote{Para soporte técnico o consultas sobre este documento.}

%%   \SetDocTitle{Mi Informe Específico}   % sobreescribe el default

\SetDocLogo{../assets/logo.pdf}
\SetDocUnit{Tecnología \& Sistemas}
\SetContactUnit{Mesa de soporte}
\SetContactMembers{%
  Equipo CoreX & Soporte & soporte@empresa.com%
}
\SetContactNote{Para soporte técnico o consultas sobre este documento.}

%%   \SetDocTitle{Mi Informe Específico}   % sobreescribe el default

\SetDocLogo{../assets/logo.pdf}
\SetDocUnit{Tecnología \& Sistemas}
\SetContactUnit{Mesa de soporte}
\SetContactMembers{%
  Equipo CoreX & Soporte & soporte@empresa.com%
}
\SetContactNote{Para soporte técnico o consultas sobre este documento.}

\SetDocType{MEMORANDO}
\SetDocTitle{Instrucción para Cierre Excepcional de Registros}
\SetDocCode{MEM-GER-001}
\SetDocArea{Gerencia General}
\SetDocAuthor{Gerencia General}
\SetDocDate{2026-04-22}

\begin{document}

\MemoBlock%
  {Jefes de Área / Supervisores — Todas las áreas}%
  {Gerencia General}%
  {Procedimiento excepcional de cierre de registros — Semana 17}%
  {MEM-GER-001}

\section{Instrucción}

Con motivo del mantenimiento programado del sistema ERP el día \textbf{viernes
2026-04-24 de 20:00 a 24:00}, se instruye lo siguiente:

\begin{enumerate}
  \item Todos los registros de producción del turno tarde del 2026-04-24 deberán
        cerrarse \textbf{antes de las 19:30} en el módulo correspondiente del sistema.
  \item Los registros que no puedan cerrarse en el sistema antes del mantenimiento
        serán capturados en el formulario físico \texttt{FRM-OPS-007} y cargados
        al sistema el lunes 2026-04-27 a primera hora.
  \item El supervisor de turno es responsable de verificar que todos los registros
        de su área estén cerrados o respaldados en papel antes de las 19:30.
  \item Cualquier incidencia deberá reportarse al Jefe de Operaciones antes de
        las 20:00 del mismo día.
\end{enumerate}

\begin{ObservationBox}[Vigencia]
  Esta instrucción aplica únicamente para el cierre del turno tarde del
  2026-04-24.  Los procedimientos normales se retoman desde el turno madrugada
  del 2026-04-25.
\end{ObservationBox}

\NoteInline{Acuse de recibo:} Confirmar recepción de este memorando respondiendo
al correo de distribución antes de las \textbf{18:00 del 2026-04-23}.

\SignatureBlock{Gerencia General}{Gerencia General}

\ContactSignature

\end{document}
